% 型10：模試型
% 時間を計って通しで解く試験冊子の形をとる．
% 単元の区切りではなく，一定期間の到達度を測る目的で使う．
% 既存の mmmockpaper / mmmockquestion / \MakeMathMockTitle を用いる．
\PassOptionsToPackage{twocolumn=false}{surikumi}
\documentclass[lualatex,paper=b5j,fontsize=10pt]{jlreq}

\usepackage{surikumi-mondaishu}

\MathMockSetup{
  name={基礎到達度テスト},
  course={第 1 回／数学 I・A},
  audience={入試基礎講座 受講者向け},
  description-label={このテストについて},
  description={時間を計って実戦形式で解き，現時点の到達度を確認する．
    採点基準と講評は次号に掲載する．解き終えた時刻を各問の欄外に記録しておくと，
    どこで時間を使ったかがあとで分かる．},
  time-label={試験時間},
  time={60 分（大問 3 題）}
}

\begin{document}

\MakeMathMockTitle

\begin{mmmockpaper}
  \begin{mmmockquestion}{1}[配点 30]
    $a$ を実数とし，$f(x)=x^2-2ax+a+2$ とおく．
    \begin{mmmockparts}
      \item $f(x)=0$ が異なる 2 個の実数解をもつような $a$ の範囲を求めよ．
      \item $f(x)=0$ の 2 解がともに正であるような $a$ の範囲を求めよ．
      \item $0\leq x\leq 2$ における $f(x)$ の最小値を $a$ で表せ．
    \end{mmmockparts}
  \end{mmmockquestion}

  \begin{mmmockquestion}{2}[配点 35]
    1 個のさいころを 3 回投げ，出た目を順に $a$，$b$，$c$ とする．
    \begin{mmmockparts}
      \item $a<b<c$ となる確率を求めよ．
      \item $a$，$b$，$c$ がこの順に等差数列をなす確率を求めよ．
      \item $abc$ が偶数となる確率を求めよ．
    \end{mmmockparts}
  \end{mmmockquestion}

  \begin{mmmockquestion}{3}[配点 35]
    $\triangle\mathrm{ABC}$ において $\mathrm{AB}=6$，$\mathrm{AC}=4$，
    $\angle\mathrm{BAC}=60^\circ$ とする．
    $\angle\mathrm{BAC}$ の二等分線と辺 $\mathrm{BC}$ の交点を
    $\mathrm{D}$ とする．
    \begin{mmmockparts}
      \item 辺 $\mathrm{BC}$ の長さを求めよ．
      \item $\mathrm{BD}:\mathrm{DC}$ を求めよ．
      \item 線分 $\mathrm{AD}$ の長さを求めよ．
    \end{mmmockparts}
  \end{mmmockquestion}
\end{mmmockpaper}

\clearpage

\MSSpreadHeading{解答上の注意}

\MSCheckItem{答案は各問ごとに区切って書く．}
\MSCheckItem{場合分けをしたときは，どの場合の話かを冒頭に書く．}
\MSCheckItem{図を描いた場合は，答案の一部として残す．}
\MSCheckItem{最後に，答が問題の条件を満たすかどうかを確かめる．}

\MSSpreadHeading{自己採点欄}

\begin{msrubric}[配点]
  \MMMockNumber{1}\quad 二次方程式の解の配置 & 30 \\
  \MMMockNumber{2}\quad 確率 & 35 \\
  \MMMockNumber{3}\quad 図形と計量 & 35 \\
  合計 & 100 \\
\end{msrubric}

\MMNote{目標は 60 点である．60 点に届かなかった大問は，該当する単元の講義記事に戻ること．}

\end{document}
